\documentclass[runningheads]{llncs}

\usepackage[T1]{fontenc}
\usepackage[utf8]{inputenc}
\usepackage{graphicx}
\usepackage{amsmath,amssymb,amsfonts}
\usepackage{booktabs}
\usepackage{array}
\usepackage{float}
\usepackage{xcolor}
\usepackage{hyperref}

\hypersetup{
    colorlinks=true,
    linkcolor=blue!60!black,
    citecolor=blue!60!black,
    urlcolor=blue!60!black
}
\renewcommand\UrlFont{\color{blue}\rmfamily}
\urlstyle{rm}

\begin{document}

\title{Sentinela Clandestino: Detec\c{c}\~{a}o de Aliciamento de Apostas Ilegais e Mitiga\c{c}\~{a}o de Espelhos Mutantes na Borda via SLMs e Grafos de Amea\c{c}as}
\titlerunning{Sentinela Clandestino: Vigilante On-Device de Espelhos e APKs}

\author{Jos\'{e} Alb\'{e}rico Rocha da Silva\inst{1}}
\authorrunning{J. A. R. da Silva}

\institute{Universidade Federal Rural de Pernambuco (UFRPE)\\
Departamento de Computa\c{c}\~{a}o -- T\'{o}picos Avan\c{c}ados em Intelig\^{e}ncia Artificial (2026.2)\\
Recife - PE, Brasil\\
\email{jose.alberico@ufrpe.br}\\
\url{https://github.com/albericolx/Topicos_Avancados_IA}}

\maketitle

\begin{abstract}
Este relat\'{o}rio t\'{e}cnico formaliza a concep\c{c}\~{a}o, fundamenta\c{c}\~{a}o bibliogr\'{a}fica e an\'{a}lise comparativa do projeto \textbf{Sentinela Clandestino (Vigilante de Espelhos e APKs)}, correspondente \`{a}s etapas de Brainwriting H\'{i}brido (\texttt{brw}) e Competidores e Contribui\c{c}\~{a}o (\texttt{cmp}) da Segunda Verifica\c{c}\~{a}o de Aprendizagem (VA2) na disciplina T\'{o}picos Avan\c{c}ados em Intelig\^{e}ncia Artificial (UFRPE 2026.2). Diante da crise cr\^{o}nica de endividamento financeiro e prolifera\c{c}\~{a}o de plataformas de apostas ilegais que burlam bloqueios governamentais via dom\'{i}nios espelhos mutantes e distribui\c{c}\~{a}o clandestina de APKs nas redes sociais, prop\~{o}e-se uma arquitetura de prote\c{c}\~{a}o reativa na borda (\textit{on-device}). O sistema intercepta tentativas de redirecionamento para URLs externas e transfer\^{e}ncias de execut\'{a}veis no momento do clique, empregando um Modelo de Linguagem de Pequeno Porte (SLM) quantizado localmente para an\'{a}lise sem\^{a}ntica da p\'{a}gina de pouso e inspe\c{c}\~{a}o de manifestos de aplica\c{c}\~{a}o, orquestrado por agentes locais via \textit{Model Context Protocol} (MCP) que consultam resumos de grafos de DNS passivo. A contribui\c{c}\~{a}o \'{e} rigorosamente delimitada frente a dez ferramentas comerciais e dez artigos cient\'{i}ficos de ponta, assumindo honestamente perdas operacionais categ\'{o}ricas em lat\^{e}ncia de redirecionamento ($100\text{--}300\text{ms}$), consumo cont\'{i}nuo de mem\'{o}ria e escala de produ\c{c}\~{a}o frente a filtros est\'{a}ticos leves.

\keywords{Apostas Clandestinas \and Engenharia Social \and Small Language Models \and On-Device AI \and Grafos de DNS Passivo \and Seguran\c{c}a Mobile \and AI for Good.}
\end{abstract}

\section{Introdu\c{c}\~{a}o}
\label{sec:introducao}

\subsection{Contextualiza\c{c}\~{a}o da Crise Socioecon\^{o}mica e Dados Oficiais}

A prolifera\c{c}\~{a}o desregulada de plataformas de apostas digitais (\textit{bets}) estabeleceu uma das mais agudas crises socioecon\^{o}micas e de vulnerabilidade patrimonial registradas na hist\'{o}ria recente do Brasil. Dados oficiais consolidados pelo Banco Central do Brasil (BACEN) em agosto de 2024 revelam que o volume financeiro transacionado por pessoas f\'{i}sicas em dire\c{c}\~{a}o a plataformas de apostas atingiu a marca de \textbf{R\$~20,8 bilh\~{o}es por m\^{e}s} exclusivamente via sistema Pix. Desse total astron\^{o}mico, relat\'{o}rios t\'{e}cnicos da Secretaria de Pr\^{e}mios e Apostas do Minist\'{e}rio da Fazenda (SPA/MF) e de \'{o}rg\~{a}os de intelig\^{e}ncia financeira apontam que o \textbf{mercado clandestino e n\~{a}o autorizado absorve entre 30\% e 50\% dos dep\'{o}sitos}, canalizando mensalmente de \textbf{R\$~6 a R\$~10 bilh\~{o}es} para organiza\c{c}\~{o}es que operam sem qualquer lastro jur\'{i}dico, sede no pa\'{i}s ou garantias de solv\^{e}ncia financeira.

Embora o Minist\'{e}rio da Fazenda e a Ag\^{e}ncia Nacional de Telecomunica\c{c}\~{o}es (Anatel) tenham notificado, em outubro de 2024, mais de 2.000 dom\'{i}nios para suspens\~{a}o compuls\'{o}ria de tr\'{a}fego junto aos provedores de conex\~{a}o, a a\c{c}\~{a}o estatal revelou-se insuficiente para conter a capilaridade do fen\^{o}meno. Os operadores clandestinos responderam com uma dispers\~{a}o agressiva de hiperlinks em plataformas de mensageria instant\^{a}nea (Telegram e WhatsApp) e feeds de redes sociais (Instagram e TikTok). Esses vetores atraem o p\'{u}blico vulner\'{a}vel para \textit{landing pages} ef\^{e}meras e estimulam a instala\c{c}\~{a}o manual (\textit{side-loading}) de pacotes de aplica\c{c}\~{a}o Android (\textit{APKs}) piratas fora das lojas oficiais. Esse ecossistema combina a captura da renda das fam\'{i}lias com fraudes deliberadas de \textbf{reten\c{c}\~{a}o il\'{i}cita de saques} (\textit{honeypots} que bloqueiam o resgate dos pr\^{e}mios) e infec\c{c}\~{a}o de terminais m\'{o}veis por \textit{trojans} banc\'{a}rios camuflados~\cite{li2025automated}.

\subsection{Insufici\^{e}ncia Estrutural das Abordagens N\~{a}o-IA}

A conten\c{c}\~{a}o dessa modalidade de fraude cibern\'{e}tica falha de maneira sistem\'{a}tica quando abordada por contramedidas computacionais convencionais desprovidas de intelig\^{e}ncia sem\^{a}ntica:
\begin{enumerate}
    \item \textbf{Inocuidade de Listas Est\'{a}ticas e Bloqueadores de Teclado:} Operadores criminosos automatizam o registro cont\'{i}nuo de espelhos mutantes por meio de Algoritmos de Gera\c{c}\~{a}o de Dom\'{i}nios (DGA), sob termina\c{c}\~{o}es baratas e permissivas (\texttt{.xyz}, \texttt{.top}). Listas est\'{a}ticas de bloqueio possuem lat\^{e}ncia temporal de horas a dias at\'{e} sua atualiza\c{c}\~{a}o, tornando-se ineficazes contra endere\c{c}os ativados minutos antes do disparo de uma campanha. \textit{Resolu\c{c}\~{a}o por IA:} Um agente local orquestrado via \textit{Model Context Protocol} (MCP) consulta resumos topol\'{o}gicos de grafos de DNS passivo (pDNS)~\cite{nabeel2020following}, calculando a proximidade estrutural do novo dom\'{i}nio com clusters criminosos j\'{a} catalogados, permitindo a classifica\c{c}\~{a}o de dia zero sem lista pr\'{e}via.
    \item \textbf{Imprevisibilidade do Pix Din\^{a}mico:} Os dep\'{o}sitos n\~{a}o s\~{a}o destinados a chaves Pix est\'{a}ticas previamente marcadas, mas intermediados por APIs de \textit{gateways} de pagamento e Provedores de Servi\c{c}os de Pagamento (PSPs) que emitem c\'{o}digos EMV din\^{a}micos associados a contas banc\'{a}rias transit\'{o}rias de laranjas. \textit{Resolu\c{c}\~{a}o por IA:} Um Modelo de Linguagem de Pequeno Porte (\textit{Small Language Model} -- SLM) quantizado~\cite{liu2024mobilellm} processa o conte\'{u}do da \textit{landing page} no pr\'{o}prio aparelho, contrastando as promessas ilus\'{o}rias de ganho com cl\'{a}usulas contratuais abusivas de reten\c{c}\~{a}o de saque (\textit{rollover} de 50x e exig\^{e}ncia de taxas pr\'{e}vias para libera\c{c}\~{a}o do saldo), alertando a v\'{i}tima sobre a fraude contratual iminente.
    \item \textbf{Polimorfismo de APKs Recompilados em Massa:} O uso de f\'{a}bricas automatizadas de malware (\textit{Automated Mass Malware Factory} -- MAMF)~\cite{li2025automated} gera bin\'{a}rios Android com assinaturas criptogr\'{a}ficas SHA-256 distintas a cada download, acoplando t\'{e}cnicas de evas\~{a}o adversarial (\textit{piggybacking}) sobre aplicativos benignos~\cite{zhang2025fighting}. \textit{Resolu\c{c}\~{a}o por IA:} O modelo neural local audita semanticamente o arquivo \texttt{AndroidManifest.xml} e os fluxos de execu\c{c}\~{a}o Dex, identificando anomalias cr\'{i}ticas de seguran\c{c}a (como requisi\c{c}\~{o}es de permiss\~{a}o de Acessibilidade para sobreposi\c{c}\~{a}o de telas banc\'{a}rias e leitura de c\'{o}digos 2FA via SMS), independentemente da altera\c{c}\~{a}o sint\'{a}tica do c\'{o}digo compilado.
\end{enumerate}

\subsection{Alinhamento ao Interesse Econ\^{o}mico do Usu\'{a}rio (Paradigma Waze)}

Historicamente, ferramentas de seguran\c{c}a voltadas \`{a} autoexclus\~{a}o de jogos (e.g., Gamban) sustentam-se em bloqueios paternalistas e coercitivos. Diante de restri\c{c}\~{o}es mec\^{a}nicas r\'{i}gidas, o usu\'{a}rio compulsivo ou seduzido por promessas de lucro tende a burlar ou desinstalar sumariamente a prote\c{c}\~{a}o. Em contrapartida, o projeto Sentinela Clandestino fundamenta-se nos preceitos da Economia Comportamental, operando de maneira an\'{a}loga a um aplicativo de navega\c{c}\~{a}o colaborativa (Waze): o condutor reduz a velocidade n\~{a}o por moralismo, mas porque \'{e} informado pontualmente de que h\'{a} um radar ou buraco \`{a} frente. 

Canaliza-se, assim, o \textbf{interesse ego\'{i}sta e a avers\~{a}o \`{a} perda do pr\'{o}prio apostador}. Ao interceptar um redirecionamento fraudulento, o sistema n\~{a}o censura o jogador, mas exibe um \textbf{Painel de Evid\^{e}ncias de Risco} demonstrando fatos objetivos: (i) que as regras ocultas daquela plataforma impossibilitam matematicamente o saque dos pr\^{e}mios; (ii) que o bin\'{a}rio ofertado cont\'{e}m instru\c{c}\~{o}es de captura de senhas banc\'{a}rias; e (iii) que a raz\~{a}o social do PSP recebedor do Pix responde a processos judiciais por estelionato e opera contas de laranjas. Ao constatar que est\'{a} sendo lesado, o pr\'{o}prio usu\'{a}rio opta por abortar a transfer\^{e}ncia.

\subsection{Contribui\c{c}\~{o}es Deste Trabalho}

Frente a esse cen\'{a}rio, este artigo apresenta o \textbf{Sentinela Clandestino}, uma arquitetura de intelig\^{e}ncia artificial na borda concebida para atuar na intercepta\c{c}\~{a}o de eventos de redirecionamento e transfer\^{e}ncia de arquivos no smartphone. As principais contribui\c{c}\~{o}es formais deste trabalho s\~{a}o:
\begin{enumerate}
    \item \textbf{Formaliza\c{c}\~{a}o do Modelo de Amea\c{c}a}: Modela-se rigorosamente a din\^{a}mica advers\'{a}ria de tripla evas\~{a}o empregada por redes clandestinas (DGA, Pix din\^{a}mico em cascata e polimorfismo MAMF), definindo os limites operacionais de isolamento do sistema operacional m\'{o}vel;
    \item \textbf{Arquitetura H\'{i}brida Fast-Path/Deep-Path na Borda}: Proposi\c{c}\~{a}o de um pipeline em dois est\'{a}gios que acopla um modelo de decis\~{a}o Sistema~1 baseado em pesos abertos (\textit{LAYA}~\cite{convai2026laya}, derivado de ModernBERT~\cite{warner2024modernbert}) para infer\^{e}ncia r\'{a}pida ($< 35\text{ms}$) a um SLM quantizado em INT4 (1B--3B)~\cite{liu2024mobilellm} ativado sob demanda, garantindo $100\%$ de privacidade local sem telemetria em nuvem;
    \item \textbf{Orquestra\c{c}\~{a}o Local via MCP com Grafos pDNS}: Estrutura\c{c}\~{a}o de um agente MCP local capaz de consultar resumos compactos de grafos de DNS passivo baseados em filtros de Bloom, viabilizando correla\c{c}\~{a}o topol\'{o}gica de infraestruturas criminosas em tempo sub-milissegundo;
    \item \textbf{An\'{a}lise Sistem\'{a}tica de Trade-offs e Fronteira Cient\'{i}fica}: Mapeamento emp\'{i}rico comparativo com 10 ferramentas industriais e 10 artigos acad\^{e}micos de ponta, assumindo explicitamente as limita\c{c}\~{o}es de maturidade e impacto computacional inerentes ao paradigma de computa\c{c}\~{a}o em borda.
\end{enumerate}
\section{Trabalhos Relacionados e Matriz de Contribui\c{c}\~{a}o}
\label{sec:trabalhos_relacionados}

\subsection{Taxonomia das Abordagens Existentes}

A conten\c{c}\~{a}o de fraudes digitais e plataformas predat\'{o}rias de apostas tem sido abordada pela literatura e pela ind\'{u}stria sob quatro paradigmas tecnol\'{o}gicos predominantes:
\begin{enumerate}
    \item \textbf{Filtros de Rede e VPN Loopback Est\'{a}tica:} Softwares como Gamban, BetBlocker, NextDNS e AdGuard interceptam pacotes na camada de transporte ou atuam na resolu\c{c}\~{a}o DNS local. Embora apresentem custo computacional negligenci\'{a}vel e operem sem violar a privacidade do usu\'{a}rio, dependem inteiramente de tabelas pr\'{e}-compiladas de dom\'{i}nios conhecidos~\cite{zhauniarovich2018survey}. Consequentemente, s\~{a}o incapazes de conter dom\'{i}nios espelhos mutantes de dia zero gerados por DGA ou aplicativos distribu\'{i}dos fora de reposit\'{o}rios centrais;
    \item \textbf{Plataformas Forenses e Auditoria Est\'{a}tica de APKs:} Solu\c{c}\~{o}es como Koodous, Exodus Privacy e VirusTotal Mobile fornecem diagn\'{o}sticos detalhados sobre permiss\~{o}es abusivas e rastreadores embutidos. Todavia, foram desenhadas para analistas de seguran\c{c}a em ambientes laboratoriais (\textit{post-mortem}), exigindo submiss\~{a}o manual ou telemetria em nuvem, sendo cegas \`{a} sem\^{a}ntica contratual da plataforma de aposta e inaptas a interceptar a v\'{i}tima no ato do clique;
    \item \textbf{Detectores Heur\'{i}sticos e Agentes em Nuvem:} Ferramentas como Netcraft Mobile, Trend Micro Check e ScamAdviser monitoram URLs e comunica\c{c}\~{o}es para conter golpes de \textit{phishing}. Embora o Trend Micro Check utilize agentes conversacionais para auditar prints e mensagens, sua opera\c{c}\~{a}o exige o envio de dados \`{a} infraestrutura centralizada da empresa, n\~{a}o realiza intercepta\c{c}\~{a}o aut\^{o}noma no fluxo de navega\c{c}\~{a}o e carece de especializa\c{c}\~{a}o na engenharia social e nos termos de reten\c{c}\~{a}o de apostas vigentes no Brasil;
    \item \textbf{Modelos Acad\^{e}micos Especializados (NLP, Grafos e SLMs):} Pesquisas recentes exploram aprendizado profundo para an\'{a}lise textual de pe\c{c}as publicit\'{a}rias~\cite{muzakir2025text}, corpora de persuas\~{a}o de apostas~\cite{alvarezojeda2025precom}, minera\c{c}\~{a}o topol\'{o}gica de dom\'{i}nios via DNS passivo~\cite{nabeel2020following,khalil2016discovering} e detectores adversariais de malware~\cite{li2025automated,zhang2025fighting}. Contudo, essas propostas permanecem fragmentadas: ou s\~{a}o desenhadas para modera\c{c}\~{a}o massiva na infraestrutura de grandes provedores, ou dependem de aglomerados computacionais incompat\'{i}veis com o processamento restrito na borda m\'{o}vel.
\end{enumerate}

\subsection{Sele\c{c}\~{a}o dos Trabalhos Mais Pr\'{o}ximos}

A partir do diagn\'{o}stico sistem\'{a}tico de 10 ferramentas e 10 artigos consolidados na literatura recente, selecionaram-se tr\^{e}s ferramentas industriais e tr\^{e}s refer\^{e}ncias cient\'{i}ficas de ponta que delimitam as fronteiras operacionais do problema (Tabela~\ref{tab:matriz_concorrentes}):
\begin{itemize}
    \item \textbf{Trend Micro Check:} Mais pr\'{o}ximo na intera\c{c}\~{a}o direta com o usu\'{a}rio m\'{o}vel contra golpes digitais, por\'{e}m requer iniciativa manual da v\'{i}tima, depende de processamento em nuvem externa e n\~{a}o intercepta \textit{side-loading} de APKs em tempo real;
    \item \textbf{Gamban:} Principal refer\^{e}ncia comercial no bloqueio estrito de plataformas de apostas no smartphone, mas restrito a listas est\'{a}ticas de dom\'{i}nios e incapaz de avaliar cl\'{a}usulas de fraude ou espelhos ef\^{e}meros;
    \item \textbf{Netcraft Mobile:} Solu\c{c}\~{a}o de maior proximidade na intercepta\c{c}\~{a}o de requisi\c{c}\~{o}es de rede no momento do clique, contudo desprovida de modelos de linguagem capazes de interpretar g\'{i}rias de aliciamento ou termos contratuais de pagamento;
    \item \textbf{Muzakir et al. (2025)~\cite{muzakir2025text}:} Artigo cient\'{i}fico mais convergente na detec\c{c}\~{a}o sem\^{a}ntica de an\'{u}ncios de apostas por meio de aprendizado profundo (AutoML e CNNs com 92,80\% de acur\'{a}cia sobre 15.286 senten\c{c}as), por\'{e}m concebido para modera\c{c}\~{a}o corporativa em nuvem, restrito \`{a} l\'{i}ngua indon\'{e}sia e dissociado do terminal m\'{o}vel;
    \item \textbf{\'{A}lvarez-Ojeda et al. (2025)~\cite{alvarezojeda2025precom}:} Trabalho seminal na caracteriza\c{c}\~{a}o lingu\'{i}stica de comunidades vulner\'{a}veis e persuas\~{a}o de apostas nas redes sociais, por\'{e}m puramente descritivo e sem qualquer componente execut\'{a}vel de interven\c{c}\~{a}o t\'{e}cnica no dispositivo;
    \item \textbf{Nabeel et al. (2020)~\cite{nabeel2020following}:} Refer\^{e}ncia can\^{o}nica em infer\^{e}ncia topol\'{o}gica sobre grafos de DNS passivo para neutraliza\c{c}\~{a}o de campanhas maliciosas de dia zero, cujos algoritmos dependem de fluxos volumosos de telemetria de operadoras de telecomunica\c{c}\~{o}es e n\~{a}o operam sob o contexto local de um cliente m\'{o}vel.
\end{itemize}

\subsection{Matriz Comparativa de Eixos e Evid\^{e}ncias Emp\'{i}ricas}

A Tabela~\ref{tab:matriz_concorrentes} formaliza a matriz comparativa de oito eixos t\'{e}cnicos derivados das tens\~{o}es arquiteturais do campo. Evidenciando a honestidade acad\^{e}mica exigida pelo m\'{e}todo cient\'{i}fico, o Sentinela Clandestino assume perdas e empates nos tr\^{e}s eixos finais de engenharia.

\begin{table*}[t]
\centering
\caption{Matriz comparativa de competidores e contribui\c{c}\~{a}o ($\checkmark$: implementa plenamente; $\blacktriangleright$: implementa parcialmente; $\times$: n\~{a}o implementa; ?: aus\^{e}ncia de evid\^{e}ncia emp\'{i}rica). Os tr\^{e}s eixos \`{a} direita evidenciam os trade-offs assumidos pela proposta.}
\label{tab:matriz_concorrentes}
\setlength{\tabcolsep}{4.5pt}
\scriptsize
\begin{tabular}{lcccccccc}
\toprule
\textbf{Solu\c{c}\~{a}o / Proposta} & \textbf{Eixo 1} & \textbf{Eixo 2} & \textbf{Eixo 3} & \textbf{Eixo 4} & \textbf{Eixo 5} & \textbf{Eixo 6} & \textbf{Eixo 7} & \textbf{Eixo 8} \\
& Intercep. & Audita & Espelhos & Cl\'{a}usulas & Dados & Lat\^{e}ncia & Baixo Cons. & Prontid\~{a}o \\
& Redirec. & APK & Grafos & Calote/Saque & Celular & $< 50\text{ms}$ & Bat/RAM & Escala Real \\
\midrule
Trend Micro Check & $\times$ & $\times$ & $\blacktriangleright$ & $\blacktriangleright$ & $\times$ & $\checkmark$ & $\checkmark$ & $\checkmark$ \\
Gamban & $\times$ & $\blacktriangleright$ & $\times$ & $\times$ & $\checkmark$ & $\checkmark$ & $\checkmark$ & $\checkmark$ \\
Netcraft Mobile & $\blacktriangleright$ & $\times$ & $\blacktriangleright$ & $\times$ & $\times$ & $\checkmark$ & $\checkmark$ & $\checkmark$ \\
Muzakir et al. (2025)~\cite{muzakir2025text} & $\times$ & $\times$ & $\times$ & $\checkmark$ & $\times$ & ? & ? & $\times$ \\
\'{A}lvarez-Ojeda et al. (2025)~\cite{alvarezojeda2025precom} & $\times$ & $\times$ & $\times$ & $\checkmark$ & $\times$ & ? & ? & $\times$ \\
Nabeel et al. (2020)~\cite{nabeel2020following} & $\times$ & $\times$ & $\checkmark$ & $\times$ & $\times$ & ? & ? & $\times$ \\
\midrule
\textbf{Sentinela Clandestino} & $\mathbf{\checkmark}$ & $\mathbf{\checkmark}$ & $\mathbf{\checkmark}$ & $\mathbf{\checkmark}$ & $\mathbf{\checkmark}$ & $\mathbf{\checkmark}$ & $\mathbf{\blacktriangleright}$ & $\mathbf{\times}$ \\
\bottomrule
\end{tabular}
\end{table*}

\noindent\textbf{Evid\^{e}ncias Emp\'{i}ricas dos Eixos:}
\begin{itemize}
    \item \textit{Trend Micro Check:} Opera avalia\c{c}\~{a}o de reputa\c{c}\~{a}o de URLs via heur\'{i}sticas em nuvem ($\blacktriangleright$), realiza triagem gen\'{e}rica de mensagens em portugu\^{e}s ($\blacktriangleright$), executa como cliente leve delegando a computa\c{c}\~{a}o a servidores externos com resposta r\'{a}pida ($\checkmark$), ultrapassando 1 milh\~{a}o de instala\c{c}\~{o}es ativas ($\checkmark$);
    \item \textit{Gamban:} Realiza bloqueio mec\^{a}nico de tr\'{a}fego de rede sem inspe\c{c}\~{a}o est\'{a}tica profunda de \textit{bytecodes} ($\blacktriangleright$), opera via interface VPN loopback sem telemetria remota de hist\'{o}rico ($\checkmark$), consumindo energia residual no processador m\'{o}vel ($\checkmark$), sendo solu\c{c}\~{a}o comercial consolidada em esquemas nacionais de autoexclus\~{a}o no Reino Unido e Europa ($\checkmark$);
    \item \textit{Netcraft Mobile:} Intercepta eventos de toque em links potencialmente perigosos ($\blacktriangleright$), verifica feeds de alta disponibilidade contra servidores de golpe conhecidos ($\blacktriangleright$), opera com impacto impercept\'{i}vel de CPU/RAM no terminal m\'{o}vel ($\checkmark$), auditado e adotado em larga escala por institui\c{c}\~{o}es financeiras globais ($\checkmark$);
    \item \textit{Muzakir et al. (2025)~\cite{muzakir2025text}:} Comprovam acur\'{a}cia de 92,80\% na extra\c{c}\~{a}o e classifica\c{c}\~{a}o sem\^{a}ntica de an\'{u}ncios de apostas com express\~{o}es de persuas\~{a}o ($\checkmark$), dependendo de infraestrutura de computa\c{c}\~{a}o intensiva desvinculada de restri\c{c}\~{o}es m\'{o}veis;
    \item \textit{\'{A}lvarez-Ojeda et al. (2025)~\cite{alvarezojeda2025precom}:} Demonstram a viabilidade de modelagem de linguagem para identificar padres persuasivos de jogos de azar em feeds digitais ($\checkmark$), mantendo car\'{a}ter puramente anal\'{i}tico;
    \item \textit{Nabeel et al. (2020)~\cite{nabeel2020following}:} Validam em base real de telecomunica\c{c}\~{o}es que algoritmos de infer\^{e}ncia em grafos de traces de DNS passivo detectam campanhas furtivas at\'{e} 14 dias antes de sua inclus\~{a}o em listas p\'{u}blicas de seguran\c{c}a ($\checkmark$).
\end{itemize}

\subsection{Declara\c{c}\~{a}o Can\^{o}nica da Contribui\c{c}\~{a}o}

A partir do mapeamento sistem\'{a}tico da literatura e da identifica\c{c}\~{a}o das falhas de cobertura dos sistemas vigentes, a contribui\c{c}\~{a}o central deste trabalho \'{e} sintetizada na seguinte declara\c{c}\~{a}o can\^{o}nica:
\begin{quote}
    \textit{``O Sentinela Clandestino correlaciona a sem\^{a}ntica da p\'{a}gina de pouso e as regras contratuais de saque com a topologia de infraestrutura de espelhos e bin\'{a}rios APK diretamente na borda do smartphone, e por isso bloqueia o ciclo da aposta clandestina no momento do redirecionamento com total privacidade local, o que nem os filtros de rede est\'{a}ticos (Gamban/Netcraft) nem os classificadores corporativos em nuvem (Muzakir/{\'A}lvarez-Ojeda) conseguem realizar isoladamente.''}
\end{quote}

O concorrente mais pr\'{o}ximo \'{e} o \textbf{Trend Micro Check}, que analisa prints e links de golpes com intelig\^{e}ncia artificial; por\'{e}m, ele falha categoricamente por exigir a\c{c}\~{a}o manual proativa da v\'{i}tima, transferir os dados confidenciais do usu\'{a}rio para a nuvem corporativa e n\~{a}o dispor de capacidade t\'{e}cnica para auditar manifestos de APKs piratas nem correlacionar espelhos mutantes no instante em que ocorrem.
\section{Arquitetura Sist\^{e}mica e Metodologia}
\label{sec:arquitetura}

A arquitetura do Sentinela Clandestino foi desenhada para operar sob as restri\c{c}\~{o}es r\'{i}gidas de seguran\c{c}a, mem\'{o}ria e consumo de energia t\'{i}picas de dispositivos m\'{o}veis populares, rejeitando regras heur\'{i}sticas cegas (como bloqueios est\'{a}ticos gen\'{e}ricos de arquivos ou listas fixas de dom\'{i}nios) e articulando tr\^{e}s componentes de IA na borda.

\subsection{Camada de Intercepta\c{c}\~{a}o e Triagem de Eventos na Borda}
Em vez de tentar violar o isolamento de processos de redes sociais propriet\'{a}rias (inviabilizado pelo \textit{sandboxing} do Android), o sistema atua na transi\c{c}\~{a}o leg\'{i}tima de contexto do sistema operacional:
\begin{enumerate}
    \item \textbf{Intercepta\c{c}\~{a}o de Redirecionamento (\textit{Intent Interception}):} Captura requisi\c{c}\~{o}es de abertura de links externos despachadas a partir de navegadores embutidos (\textit{Custom Tabs}) ou mensageiros instant\^{a}neos;
    \item \textbf{Gatilho de Download de Bin\'{a}rios:} Detecta o in\'{i}cio de transfer\^{e}ncias de arquivos execut\'{a}veis com extens\~{a}o \texttt{.apk} provenientes de fontes desconhecidas fora da loja Google Play.
\end{enumerate}

\subsection{Pipeline H\'{i}brido de Dois Est\'{a}gios: Fast-Path com LAYA e Deep-Path com SLM}
Para conciliar an\'{a}lise sem\^{a}ntica profunda com lat\^{e}ncia impercept\'{i}vel na navega\c{c}\~{a}o móvel, adota-se um pipeline h\'{i}brido executado 100\% \textit{on-device}:
\begin{enumerate}
    \item \textbf{Fast-Path com Modelo de Decis\~{a}o LAYA ($< 35\text{ms}$):} No instante do redirecionamento, um modelo n\~{a}o autorregressivo de Sistema~1 baseado em pesos abertos (\textit{LAYA}, com codificadores ModernBERT) processa o DOM e manifesto em paralelo, pontuando probabilidades estruturadas de golpe com pegada inferior a $100\text{MB}$ de RAM. Se o risco for inequ\'{i}voco ($> 90\%$), o alerta \'{e} disparado instantaneamente sem depend\^{e}ncia de nuvens propriet\'{a}rias (como JEV);
    \item \textbf{Deep-Path com SLM Quantizado (Sob Demanda):} Em casos amb\'{i}guos ou sob solicita\c{c}\~{a}o de auditoria detalhada pelo usu\'{a}rio, um SLM quantizado em 4-bit (INT4, 1B--3B)~\cite{liu2024mobilellm} \'{e} acordado para dissecar as cl\'{a}usulas contratuais de pagamento em busca de armadilhas de \textbf{reten\c{c}\~{a}o il\'{i}cita de saques} (\textit{rollovers} abusivos e cobran\c{c}a de taxas pr\'{e}vias).
\end{enumerate}

\subsection{Inspe\c{c}\~{a}o de Código Polimórfico e Correla\c{c}\~{a}o em Grafos}
Quando o evento envolve um pacote de aplica\c{c}\~{a}o Android:
\begin{enumerate}
    \item \textbf{An\'{a}lise Sem\^{a}ntica do Manifesto e Bytecode:} Como os criminosos recompilam APKs dinamicamente para alterar assinaturas SHA-256 e ofuscar chamadas de API (\textit{automated mass malware factory})~\cite{li2025automated,zhang2025fighting}, o SLM analisa a discrep\^{a}ncia funcional entre a descri\c{c}\~{a}o do app e as permiss\~{o}es cr\'{i}ticas requisitadas (acesso a SMS para intercepta\c{c}\~{a}o de 2FA, sobreposi\c{c}\~{a}o de tela banc\'{a}ria e leitura de \textit{clipboard});
    \item \textbf{Agente MCP e Consulta a Grafos de Amea\c{c}as:} Por meio de um agente orquestrado via \textit{Model Context Protocol} (MCP), o sistema consulta um \textit{digest} comprimido (baseado em filtros de Bloom e grafos de resolu\c{c}\~{a}o passiva de DNS~\cite{nabeel2020following}) para correlacionar o dom\'{i}nio do servidor com clusters hist\'{o}ricos de operadores clandestinos e laranjas financeiros.
\end{enumerate}

Em vez de impor um bloqueio opaco que gere frustra\c{c}\~{a}o e leve \`{a} desinstala\c{c}\~{a}o do sistema, o Sentinela Clandestino apresenta um **Painel de Evid\^{e}ncias de Risco**: comprova ao apostador se a plataforma possui ind\'{i}cios de n\~{a}o pagamento de saques, c\'{o}digo espi\~{a}o ou chave Pix vinculada a contas sob investiga\c{c}\~{a}o de lavagem de dinheiro, canalizando o interesse do pr\'{o}prio usu\'{a}rio em n\~{a}o ser v\'{i}tima de estelionato.

\section{Considera\c{c}\~{o}es Finais e Trade-offs de Engenharia}
\label{sec:conclusao}

A viabiliza\c{c}\~{a}o do Sentinela Clandestino imp\~{o}e a aceita\c{c}\~{a}o expl\'{i}cita de compromissos t\'{e}cnicos de engenharia de software e sistemas operacionais, evitando concess\~{o}es complacentes:

\begin{enumerate}
    \item \textbf{Lat\^{e}ncia de Redirecionamento ($100\text{--}300\text{ms}$ vs. $< 50\text{ms}$):} Filtros est\'{a}ticos de DNS ou VPN (como Gamban e Netcraft) operam com descarte quase instant\^{a}neo de pacotes ($< 10\text{ms}$). A execu\c{c}\~{a}o de uma passagem direta (\textit{forward pass}) de um SLM quantizado e a consulta ao grafo local introduzem uma lat\^{e}ncia percept\'{i}vel de at\'{e} 300~milissegundos durante o clique externo, configurando um custo deliberado em prol da an\'{a}lise contextual profunda;
    \item \textbf{Consumo Energ\'{e}tico e Aloca\c{c}\~{a}o de Mem\'{o}ria:} Manter um modelo de linguagem residente na mem\'{o}ria de smartphones de entrada causaria o encerramento do processo pelo \textit{Low Memory Killer} (LMK) do Android. Portanto, a arquitetura adota ativa\c{c}\~{a}o estritamente orientada a eventos, permanecendo inativa durante a navega\c{c}\~{a}o cotidiana e ativando o processador neural (NPU/CPU) apenas em eventos de clique externo ou download;
    \item \textbf{Maturidade de Valida\c{c}\~{a}o em Escala:} Enquanto produtos comerciais possuem bases de milh\~{o}es de usu\'{a}rios e certifica\c{c}\~{a}o regulat\'{o}ria, o Sentinela Clandestino posiciona-se como uma investiga\c{c}\~{a}o tecnol\'{o}gica de vanguarda voltada para o bem social (\textit{AI for Good}).
\end{enumerate}

Em s\'{i}ntese, enquanto as solu\c{c}\~{o}es existentes tentam conter o avan\c{c}o das fraudes mediante tabelas fixas de endere\c{c}os ou delegando a privacidade a servidores corporativos na nuvem, o Sentinela Clandestino articula a triagem sem\^{a}ntica no aparelho com a intelig\^{e}ncia de grafos de infraestrutura criminosa, oferecendo prote\c{c}\~{a}o em tempo real sem comprometer a soberania dos dados do cidad\~{a}o.


\bibliographystyle{splncs04}
\bibliography{references}

\end{document}
